\chapter*{前言}

第二册原理卷负责回答“为什么”：为什么热循环会被分支、依赖和内存等待限制，为什么 false sharing 会让正确程序扩展失败，为什么条件变量必须围绕等待谓词，为什么异步 I/O 的 submit、completion 和 commit 不是同一件事，为什么分布式系统必须区分逻辑任务和物理 attempt。本卷负责回答“怎样落地”：怎样把这些原则写成一个可以构建、测试、测量、故障注入和复盘的 C++20 工程。

本卷的贯穿项目统一称为 \term{Compute Systems Engine}{计算系统引擎}。它从一个日志统计任务开始：输入是一批记录，输出是按 key 聚合的稳定结果。第一阶段只要求顺序 reference、输入合同和诊断；后续逐步加入 split、batch、冷热字段、benchmark、并行内核、线程池、有界队列、原子协议、任务图、可靠写、manifest、checkpoint、MessageBus、attempt、epoch、分片和故障注入。每一步都必须有测试和报告，不接受只写代码、不解释边界的实现。

本卷不是代码片段合集，而是一部工程教材。每章都遵循同一条路线：先说明本章要兑现原理卷中的哪个机制；再给出接口和目录组织；随后实现最小版本、测试失败输入、补充报告字段；最后说明这个版本还不能处理什么，下一章为什么必须继续扩展。读者应当把代码、测试和报告看成一个整体：代码实现机制，测试固定语义，报告保存证据。

\begin{keyidea}
代码实践卷的目标不是写一个庞大的教学框架，而是让每一项复杂度都有购买理由：它修复了哪个失败，保留了哪个正确性边界，带来了什么运行成本，怎样用测试和报告证明。
\end{keyidea}
